\documentclass[11pt,a4paper]{article}
\usepackage[margin=1in]{geometry}
\usepackage[T1]{fontenc}
\usepackage[utf8]{inputenc}
\usepackage{lmodern}
\usepackage{amsmath,amssymb}
\usepackage{booktabs}
\usepackage{enumitem}
\usepackage{hyperref}
\usepackage{xcolor}
\usepackage{listings}
\usepackage{graphicx}

\hypersetup{
  colorlinks=true,
  linkcolor=blue!60!black,
  urlcolor=blue!60!black
}

\lstdefinestyle{pycode}{
  language=Python,
  basicstyle=\ttfamily\small,
  keywordstyle=\color{blue!70!black},
  commentstyle=\color{green!40!black},
  stringstyle=\color{orange!50!black},
  showstringspaces=false,
  breaklines=true,
  frame=single
}
\newcommand{\img}[2][]{\includegraphics[#1]{\detokenize{#2}}}

\title{Filter Theory and Implementation Report\\\large xotiena00}
\author{FaceFiltering Project Analysis}
\date{\today}

\begin{document}
\maketitle

\section{Scope}
This report redesigns \texttt{xotiena00} for the requested filter groups:
\begin{itemize}[leftmargin=1.2em]
  \item \textbf{Smoothing / Denoising}: Bilateral, Mean, Gaussian
  \item \textbf{Tone / Brightness Adjustment}: Gamma, Dodge
\end{itemize}

The report compares and classifies filters, summarizes theory and practical application,
and discusses computational cost and implementation strategy.

\section{Classification and Comparison}
\begin{center}
\begin{tabular}{@{}lllll@{}}
\toprule
\textbf{Filter} & \textbf{Group} & \textbf{Methodology} & \textbf{Linear?} & \textbf{Edge Preservation} \\
\midrule
Mean & Smooth/Denoise & Spatial neighborhood averaging & Yes & Low \\
Gaussian & Smooth/Denoise & Spatial convolution (Gaussian kernel) & Yes & Low--Medium \\
Bilateral & Smooth/Denoise & Nonlinear spatial+range weighting & No & High \\
Gamma & Tone/Brightness & Point-wise power-law mapping (LUT) & No & N/A \\
Dodge & Tone/Brightness & Point-wise nonlinear brightening curve & No & N/A \\
\bottomrule
\end{tabular}
\end{center}

\section{Theory and Application}

\subsection{Mean Filter}
\textbf{Theory:}
\[
I'(p)=\frac{1}{|\Omega|}\sum_{q\in\Omega} I(q)
\]
It replaces each pixel by the arithmetic mean of a local window \(\Omega\).

\textbf{Pixel-level transformation (exact):}
For each output location \(p=(x,y)\), collect all pixels \(q\) in the \(k\times k\) neighborhood centered at \(p\), sum their intensities (or channel values), divide by \(k^2\), and assign:
\[
I'_{c}(x,y)=\frac{1}{k^2}\sum_{i=-r}^{r}\sum_{j=-r}^{r} I_{c}(x+i,y+j),\quad r=\frac{k-1}{2}
\]
So the center pixel is replaced by the plain local average; sharp transitions are reduced because values from both sides of an edge are mixed.

\textbf{Application:}
simple denoising and pre-smoothing before segmentation; useful when speed matters more than edge quality.

\textbf{Compare with:}
Gaussian gives visually smoother blur; bilateral preserves edges much better.

\subsection{Gaussian Filter}
\textbf{Theory:}
\[
I' = G_{\sigma} * I,\quad
G(x,y)=\frac{1}{2\pi\sigma^2}\exp\left(-\frac{x^2+y^2}{2\sigma^2}\right)
\]
It is a weighted average where nearby pixels contribute more.

\textbf{Pixel-level transformation (exact):}
For each output pixel \(p\), take a weighted sum of its neighborhood, where weights are Gaussian and normalized to sum to 1:
\[
I'_{c}(x,y)=\sum_{i=-r}^{r}\sum_{j=-r}^{r} G_{\sigma}(i,j)\,I_{c}(x+i,y+j)
\]
Pixels close to \((x,y)\) get larger weights than far pixels, so smoothing is stronger locally while still blurring edges.
\textbf{Practical note:} in many implementations, increasing kernel size \(k\) often causes a more obvious increase in blur than small changes in \(\sigma\), because a larger window includes more distant neighbors. In practice, blur appearance is controlled by both \(k\) and \(\sigma\), and they should be tuned together.

\textbf{Application:}
general denoising, anti-aliasing, and pre-processing for edge detectors.

\textbf{Compare with:}
stronger quality than mean at similar intent; still blurs edges unlike bilateral.

\subsection{Bilateral Filter}
\textbf{Theory:}
\[
I'(p)=\frac{\sum_{q\in\Omega} w_s(\|p-q\|)\,w_r(|I(p)-I(q)|)\,I(q)}
{\sum_{q\in\Omega} w_s(\|p-q\|)\,w_r(|I(p)-I(q)|)}
\]
with
\[
w_s(d)=\exp\!\left(-\frac{d^2}{2\sigma_s^2}\right),\quad
w_r(\Delta I)=\exp\!\left(-\frac{\Delta I^2}{2\sigma_r^2}\right)
\]
It combines spatial closeness and intensity similarity, so cross-edge mixing is reduced.

\textbf{Pixel-level transformation (exact):}
For each output pixel \(p\), each neighbor \(q\) gets a product weight:
\[
w(p,q)=w_s(\|p-q\|)\cdot w_r(|I(p)-I(q)|)
\]
Then the output is normalized weighted averaging:
\[
I'_{c}(p)=\frac{\sum_{q\in\Omega} w(p,q)\,I_{c}(q)}{\sum_{q\in\Omega} w(p,q)}
\]
Thus, a nearby pixel with very different intensity contributes little, which is why bilateral preserves edges better than mean/Gaussian.

\textbf{Application:}
edge-preserving denoising (portraits, object boundaries, stylized smoothing).

\textbf{Compare with:}
better edge retention than mean/gaussian; higher computational cost.
\textbf{Performance comparison note:}
\begin{itemize}[leftmargin=1.2em]
  \item Bilateral typically performs better than Gaussian for denoising when edge preservation is important.
  \item Compared with Median, Bilateral often preserves smooth gradients and fine texture better.
  \item For strong salt-and-pepper (impulse) noise, Median can outperform Bilateral.
  \item Bilateral is computationally more expensive, so Gaussian/Median may be preferred for real-time speed.
\end{itemize}

\subsection{Gamma Filter}
\textbf{Theory:}
\[
V_{\text{out}} = A\,V_{\text{in}}^{\gamma}
\]
Gamma correction is a nonlinear luminance mapping used in image/video encoding and decoding.
In the common case \(A=1\), both \(V_{\text{in}}\) and \(V_{\text{out}}\) are normalized to \([0,1]\).

\textbf{Encoding vs decoding interpretation:}
\begin{itemize}[leftmargin=1.2em]
  \item \(\gamma<1\): compressive mapping (often called \emph{gamma compression}/encoding),
  \item \(\gamma>1\): expansive mapping (often called \emph{gamma expansion}/decoding).
\end{itemize}

In this project, for 8-bit intensities \(I\in[0,255]\), the implementation uses:
\[
I' = 255\left(\frac{I}{255}\right)^{1/\gamma}
\]
This keeps the operation point-wise while matching practical display-style control.

\textbf{Behavioral intuition:}
\begin{itemize}[leftmargin=1.2em]
  \item \(\gamma<1\) brightens midtones more strongly.
  \item \(\gamma>1\) darkens midtones and preserves highlights.
  \item \(\gamma=1\) leaves the image unchanged.
\end{itemize}

\textbf{Pixel-level transformation (exact):}
Each pixel is transformed independently (no neighborhood):
\[
I'_{c}(x,y)=255\left(\frac{I_{c}(x,y)}{255}\right)^{1/\gamma}
\]
\textbf{Notation:} \(I_{c}(x,y)\) is the input pixel intensity at position \((x,y)\) for channel \(c\) (typically in \([0,255]\) for 8-bit images), and \(I'_{c}(x,y)\) is the output intensity after gamma correction.
\textbf{Why \(0..255\):} the implementation uses 8-bit image storage (\texttt{uint8}), so each channel has 256 discrete levels from 0 to 255. Therefore formulas normalize by dividing by 255 and convert back by multiplying by 255. For other bit depths, the scale changes (e.g., 16-bit uses \(0..65535\)).

Implementation-wise, this is usually done via LUT:
\[
\text{LUT}[i]=255\left(\frac{i}{255}\right)^{1/\gamma},\quad I'_{c}(x,y)=\text{LUT}[I_{c}(x,y)]
\]
So every channel value is remapped by the same nonlinear curve.

\textbf{Application:}
midtone correction, global brightening/darkening, display-oriented tone adjustment.

\textbf{Compare with:}
less aggressive than dodge; unlike smoothing filters, it does not denoise or blur.

\subsection{Dodge Filter}
\textbf{Theory (project-style mapping):}
\[
I' \approx \frac{255\,I}{255-sI}
\]
with strength \(s\in[0,1)\). Bright values are amplified more strongly than dark values.
\textbf{Dodge strength interpretation:}
\begin{itemize}[leftmargin=1.2em]
  \item \(s\) is the parameter controlling how aggressively pixels are brightened.
  \item Higher \(s\) makes the denominator \(255-sI\) smaller for bright pixels, so brightening becomes much stronger.
  \item Lower \(s\) gives a gentler effect.
  \item \(s=0\) gives almost no change.
  \item Near the upper practical limit (e.g., \(s\approx 0.95\)), highlights can quickly saturate to white (clipping).
\end{itemize}

\textbf{Pixel-level transformation (exact):}
Each pixel/channel is transformed independently by a rational brightening function:
\[
I'_{c}(x,y)=\mathrm{clip}\!\left(\frac{255\,I_{c}(x,y)}{\max(255-s\,I_{c}(x,y),1)}\right)
\]
When \(I_c(x,y)\) is high, denominator decreases and output rises rapidly toward white. When \(I_c(x,y)\) is low, change is smaller.

\textbf{Application:}
highlight enhancement, artistic glow-like brightening, sketch-like stylization.
\textbf{Historical note:} The name ``Dodge'' comes from traditional darkroom photography, where it meant making selected regions of the light-sensitive paper ``dodge'' (avoid) light from the enlarger, so those regions developed brighter.

\textbf{Compare with:}
more aggressive than gamma in highlights; unlike bloom, dodge is point-wise (no spatial blur).

\section{Comparison Section (Visual Results)}
This section maps the collected report images to the relevant filters and highlights practical behavior differences.

\subsection{Smoothing / Denoising Comparison}
\textbf{Gaussian (kernel and blur behavior):}
\begin{center}
\img[width=0.31\textwidth]{xotiena00reportimages/gaussian kernel -small sigma.jpg}
\hfill
\img[width=0.31\textwidth]{xotiena00reportimages/gaussian kernel larger sigma.jpg}
\hfill
\img[width=0.31\textwidth]{xotiena00reportimages/gaussian blur 1- kernel.jpg}

\vspace{0.5em}
\img[width=0.31\textwidth]{xotiena00reportimages/gaussian blur 1.jpg}
\hfill
\img[width=0.31\textwidth]{xotiena00reportimages/gaussian blur 2.jpg}
\hfill
\img[width=0.31\textwidth]{xotiena00reportimages/gaussian blur no smooth edges.jpg}
\end{center}

\textbf{Median and mean-context comparison (uniform smoothing vs edge/detail trade-off):}
\begin{center}
\img[width=0.31\textwidth]{xotiena00reportimages/median lower kernel.jpg}
\hfill
\img[width=0.31\textwidth]{xotiena00reportimages/median-kernel 9x9.jpg}
\hfill
\img[width=0.31\textwidth]{xotiena00reportimages/median noise removal.jpg}

\vspace{0.5em}
\img[width=0.48\textwidth]{xotiena00reportimages/gaussian and median comparision.jpg}
\hfill
\img[width=0.48\textwidth]{xotiena00reportimages/gaussian nosie removel.jpg}
\end{center}

\textbf{Bilateral (edge preservation):}
\begin{center}
\img[width=0.48\textwidth]{xotiena00reportimages/bilateral -edge preservations.jpg}
\hfill
\img[width=0.48\textwidth]{xotiena00reportimages/bilatera -edge preservations kernel.jpg}
\end{center}

\textbf{Interpretation:}
\begin{itemize}[leftmargin=1.2em]
  \item Gaussian and mean-style smoothing suppress noise but can soften boundaries.
  \item Median better handles impulse noise but may flatten small textures.
  \item Bilateral keeps stronger edges while still reducing noise in flatter regions.
\end{itemize}

\subsection{Tone / Brightness Comparison}
\textbf{Gamma progression:}
\begin{center}
\img[width=0.19\textwidth]{xotiena00reportimages/gamma1.jpg}
\hfill
\img[width=0.19\textwidth]{xotiena00reportimages/gamma2.jpg}
\hfill
\img[width=0.19\textwidth]{xotiena00reportimages/gamma3.jpg}
\hfill
\img[width=0.19\textwidth]{xotiena00reportimages/gamma5.jpg}
\hfill
\img[width=0.19\textwidth]{xotiena00reportimages/gamma6.jpg}
\end{center}

\textbf{Dodge progression and gamma-vs-dodge extremes:}
\begin{center}
\img[width=0.24\textwidth]{xotiena00reportimages/dodge1.jpg}
\hfill
\img[width=0.24\textwidth]{xotiena00reportimages/dodge2.jpg}
\hfill
\img[width=0.24\textwidth]{xotiena00reportimages/dodge3.jpg}
\hfill
\img[width=0.24\textwidth]{xotiena00reportimages/dodge-failes on eyes reveal.jpg}

\vspace{0.5em}
\img[width=0.31\textwidth]{xotiena00reportimages/dodge-gamma-lower values.jpg}
\hfill
\img[width=0.31\textwidth]{xotiena00reportimages/dodge-gamma-higher values.jpg}
\hfill
\img[width=0.31\textwidth]{xotiena00reportimages/dodge-gamma-xtremehighjpg.jpg}
\end{center}

\textbf{Interpretation:}
\begin{itemize}[leftmargin=1.2em]
  \item Gamma produces controlled global tone remapping, especially in midtones.
  \item Dodge is more aggressive in highlights and can saturate bright facial regions.
  \item Extreme dodge strength can introduce unnatural highlight clipping artifacts.
\end{itemize}

\section{Computational Power and Complexity}

\subsection{Asymptotic Cost per Pixel}
Let \(k\) be kernel width (window \(k\times k\)).
\begin{itemize}[leftmargin=1.2em]
  \item Mean: \(O(k^2)\) naive, can be reduced using integral images.
  \item Gaussian: \(O(k^2)\) for direct 2D convolution, or \(O(2k)\) with separable implementation.
  \item Bilateral: \(O(k^2)\) with expensive nonlinear weights per neighbor.
  \item Gamma: \(O(1)\) with LUT lookup.
  \item Dodge: \(O(1)\) per pixel arithmetic.
\end{itemize}

\subsection{Practical Runtime Ranking}
Typical CPU runtime (fastest to slowest):
\[
\text{Gamma} \approx \text{Dodge} \ll \text{Mean} \lesssim \text{Gaussian} \ll \text{Bilateral}
\]
Actual order depends on implementation backend (NumPy loops vs vectorized/OpenCV C++ kernels).

\section{Implementation Methodology}
\subsection{Shared pipeline}
\begin{enumerate}[leftmargin=1.5em]
  \item normalize input to BGR \texttt{uint8},
  \item clamp parameters to safe ranges,
  \item compute in floating-point where needed,
  \item clip and cast back to \([0,255]\) \texttt{uint8}.
\end{enumerate}

\subsection{Filter-specific implementation notes}
\begin{itemize}[leftmargin=1.2em]
  \item \textbf{Mean:} box kernel convolution (\texttt{1/k\^{}2} weights); good baseline.
  \item \textbf{Gaussian:} generated Gaussian kernel; prefer separable implementation for speed.
  \item \textbf{Bilateral:} requires both spatial and range Gaussian weights; use optimized library path when possible for interactive UI.
  \item \textbf{Gamma:} precompute 256-entry LUT and apply via indexing.
  \item \textbf{Dodge:} stable denominator handling and clamping are required to avoid overflow artifacts.
\end{itemize}

\section{Algorithm Summaries}

\subsection{Smoothing/Denoising}
\begin{lstlisting}[style=pycode]
# Mean
out = convolve(image, ones(k,k)/(k*k))

# Gaussian
kernel = gaussian_kernel(k, sigma)
out = convolve(image, kernel)   # or separable 1D passes

# Bilateral
for each pixel p:
    for q in kxk neighborhood:
        w = spatial_weight(p,q,sigma_space) * range_weight(I[p], I[q], sigma_color)
    out[p] = weighted_sum / weight_sum
\end{lstlisting}

\subsection{Tone/Brightness}
\begin{lstlisting}[style=pycode]
# Gamma (LUT)
lut[i] = round(((i/255.0)**(1.0/gamma))*255)
out = lut[image]

# Dodge
denom = 255 - strength*image
out = clip((image*255) / max(denom, 1))
\end{lstlisting}

\section{Conclusion}
The two groups serve different goals. Mean, Gaussian, and Bilateral are neighborhood-based
denoisers with a quality-speed trade-off where bilateral gives best edge preservation at the
highest cost. Gamma and Dodge are point-wise tone operators: gamma is controlled global tone
correction, while dodge is an aggressive highlight boost. For real-time systems, LUT-based
gamma/dodge are cheapest; for denoising quality, choose bilateral when computational budget allows,
otherwise use Gaussian as the standard compromise.

\end{document}
